% !TeX root = ../概率论与数理统计笔记.tex
% 由主文档通过 \input 载入；请编译主文档。
\setcounter{chapter}{16} % 保留原书章号。
\chapter{方差分析与试验设计初步}
\label{chap:ps-anova}

本章研究多个总体均值的比较，以及如何通过试验设计提高比较的可信度。
先修内容为正态抽样分布、假设检验与线性模型；学习目标是把模型、平方和分解和研究设计联系起来。

\section{单因素方差分析模型}
\label{sec:ps-oneway-anova}
建立\emph{单因素方差分析}（one-way analysis of variance, ANOVA）的固定效应模型。
区分组别、处理、观测单位与随机误差，明确独立性、共同误差方差及正态性用于精确检验的条件。
用基准组或效应和为零的约束处理参数表示的非唯一性。

\section{平方和分解与 F 检验}
\label{sec:ps-anova-decomposition}
将总变异分解为组间与组内变异，推导平方和、自由度和均方，构造各组均值相等的 $F$ 检验。
由线性模型正交分解解释正态误差下相关平方和的分布与独立性。
说明总体拒绝后只能断言并非所有均值都相等，不能直接确定具体哪一对不同。

\section{对比、多重比较与条件检查}
\label{sec:ps-anova-contrasts}
定义组均值的线性\emph{对比}（contrast），说明系数和为零的要求，区分事先规划与事后探索。
构造对比的估计与区间，结合第\ref{chap:ps-testing}章讨论多重比较控制。
通过残差检查讨论方差不齐、非正态和不平衡设计；替代方法需同时说明新的假设和推断目标。

\section{随机化、重复、区组与双因素设计}
\label{sec:ps-experimental-design}
介绍\emph{试验设计}（experimental design）的随机化、独立重复与\emph{区组}（block）原则。
区分技术重复与独立实验单位，说明伪重复会低估不确定性。
双因素设计中介绍主效应与交互作用；能否识别和检验交互作用取决于设计及重复情况，不能只看表格形式。

\section{例题、习题与学习检查}
规划三个处理组的均值比较、随机区组对照、存在交互作用的双因素示例。
习题要求先写模型再列方差分析表，核对各自由度之和，并说明实验单位及随机分组机制。
学习检查要求把显著性、差异大小、区间与设计限制一并解释。





